\documentclass{article}
\usepackage[utf8]{inputenc}
\usepackage{amsmath,amssymb,amsfonts,amsthm,geometry}
\geometry{margin=1in}
\title{Right Shift Operator on l^2}
\author{MathStudio Agentic Researcher}
\date{\today}
\begin{document}
\maketitle
# Definition and Properties of the Right Shift Operator on $\ell^2$

The right shift operator is a fundamental example in operator theory, particularly in the study of non-normal operators on Hilbert spaces.

## 1. Formal Definition
Let $H = \ell^2(\mathbb{N})$ (or $\ell^2(\mathbb{Z}_+)$) be the Hilbert space of square-summable sequences of complex numbers. The **right shift operator** $R: \ell^2 \to \ell^2$ is defined by:
$$R(x_1, x_2, x_3, \dots) = (0, x_1, x_2, x_3, \dots)$$
for any sequence $x = (x_n)_{n=1}^\infty \in \ell^2$.

## 2. Notation and Adjoint
- **Notation**: Commonly denoted by $R$ or $S$.
- **Adjoint**: The adjoint of the right shift operator is the **left shift operator**, denoted by $L$ or $R^*$, defined as:
  $$R^*(x_1, x_2, x_3, \dots) = (x_2, x_3, x_4, \dots)$$
  This relationship is verified by the inner product: $\langle Rx, y \rangle = \langle x, R^*y \rangle$.

## 3. Basic Properties
- **Isometry**: $R$ is an isometry, meaning $\|Rx\| = \|x\|$ for all $x \in \ell^2$.
- **Non-Normality**: 
  - $R^*R = I$ (the identity operator).
  - $RR^* \neq I$. Specifically, $RR^*(x_1, x_2, \dots) = (0, x_2, x_3, \dots)$, which is the projection onto the subspace of sequences with the first term equal to zero.
  - Since $R^*R \neq RR^*$, the operator is **not normal** and therefore **not unitary**.
- **Spectrum**:
  - The spectrum of $R$ is the closed unit disk: $\sigma(R) = \{ \lambda \in \mathbb{C} : |\lambda| \le 1 \}$.
  - $R$ has **no eigenvalues**. If $Rx = \lambda x$, then $(0, x_1, x_2, \dots) = (\lambda x_1, \lambda x_2, \dots)$, which implies $x_n = 0$ for all $n$ if $\lambda \neq 0$, and $x=0$ if $\lambda=0$.
- **Range**: The range $\text{ran}(R)$ is a closed subspace of codimension 1, specifically $\{x \in \ell^2 : x_1 = 0\}$.

## Sources
- **Barry Simon**, *Operator Theory: A Comprehensive Course in Analysis, Part 4*, p. 68 [KB ID: 19037].
- **Theo Bühler & Dietmar A. Salamon**, *Functional Analysis*, p. 242 [KB ID: 21250].

\end{document}